% Options for packages loaded elsewhere
\PassOptionsToPackage{unicode}{hyperref}
\PassOptionsToPackage{hyphens}{url}
\documentclass[
  11pt,
  letterpaper,
]{article}
\usepackage{xcolor}
\usepackage[margin=1in]{geometry}
\usepackage{amsmath,amssymb}
\setcounter{secnumdepth}{-\maxdimen} % remove section numbering
\usepackage{iftex}
\ifPDFTeX
  \usepackage[T1]{fontenc}
  \usepackage[utf8]{inputenc}
  \usepackage{textcomp} % provide euro and other symbols
\else % if luatex or xetex
  \usepackage{unicode-math} % this also loads fontspec
  \defaultfontfeatures{Scale=MatchLowercase}
  \defaultfontfeatures[\rmfamily]{Ligatures=TeX,Scale=1}
\fi
\usepackage{lmodern}
\ifPDFTeX\else
  % xetex/luatex font selection
\fi
% Use upquote if available, for straight quotes in verbatim environments
\IfFileExists{upquote.sty}{\usepackage{upquote}}{}
\IfFileExists{microtype.sty}{% use microtype if available
  \usepackage[]{microtype}
  \UseMicrotypeSet[protrusion]{basicmath} % disable protrusion for tt fonts
}{}
\makeatletter
\@ifundefined{KOMAClassName}{% if non-KOMA class
  \IfFileExists{parskip.sty}{%
    \usepackage{parskip}
  }{% else
    \setlength{\parindent}{0pt}
    \setlength{\parskip}{6pt plus 2pt minus 1pt}}
}{% if KOMA class
  \KOMAoptions{parskip=half}}
\makeatother
\usepackage{longtable,booktabs,array}
\usepackage{calc} % for calculating minipage widths
% Correct order of tables after \paragraph or \subparagraph
\usepackage{etoolbox}
\makeatletter
\patchcmd\longtable{\par}{\if@noskipsec\mbox{}\fi\par}{}{}
\makeatother
% Allow footnotes in longtable head/foot
\IfFileExists{footnotehyper.sty}{\usepackage{footnotehyper}}{\usepackage{footnote}}
\makesavenoteenv{longtable}
\setlength{\emergencystretch}{3em} % prevent overfull lines
\providecommand{\tightlist}{%
  \setlength{\itemsep}{0pt}\setlength{\parskip}{0pt}}
\usepackage{bookmark}
\IfFileExists{xurl.sty}{\usepackage{xurl}}{} % add URL line breaks if available
\urlstyle{same}
\hypersetup{
  hidelinks,
  pdfcreator={LaTeX via pandoc}}

\author{}
\date{}

\begin{document}

\section{Transformer Architectures as Tholonic Instantiations: A Formal
Role Assignment and Structural
Derivation}\label{transformer-architectures-as-tholonic-instantiations-a-formal-role-assignment-and-structural-derivation}

\textbf{Author:} J. W. Milton, Clarity Coalition

\textbf{Version:} 1.0

\textbf{Date:} 27 June 2026

\textbf{Keywords:} tholonic model, N-D-C triad, transformer
architecture, role assignment, structural derivation, virial balance,
five constants, golden ratio, phase boundaries, alignment, C-dominance,
LayerNorm, attention, MLP

\begin{center}\rule{0.5\linewidth}{0.5pt}\end{center}

\subsection{Abstract}\label{abstract}

The tholonic framework
(\href{https://github.com/baardev/truevalue/raw/main/docnav/Research/papers/1_recursive-tholonic-five-constants/1_recursive-tholonic-five-constants.pdf}{Papers
1} and
\href{https://github.com/baardev/truevalue/raw/main/docnav/Research/papers/3_minimal-recursive-triadic-framework/3_minimal-recursive-triadic-framework.pdf}{3}
in this series) establishes that any stable self-sustaining recursive
system must instantiate three functionally distinct roles: a
\emph{negotiation} state \(N\) that is iteratively refined, a
\emph{definition/limitation} variable \(D\) that bounds, and a
\emph{contribution/integration} variable \(C\) that accumulates. The
framework further proves, via convergence analysis, that the fixed
points of N-D-C recurrences are the five classical constants
\(\varphi\), \(\sqrt{2}\), \(\ln 2\), \(e\), and \(\pi/4\).

\href{https://github.com/baardev/truevalue/raw/main/docnav/Research/papers/10_tholonic-neural-architecture/10_tholonic-neural-architecture.pdf}{Paper
10} in this series reports that empirical phase boundaries in
transformer inference match one of these five constants at a 75.5\% rate
across 20 models spanning nine architecture families. That result is
presented as a test of the tholonic hypothesis applied to neural
networks. What has not been established, and what the present paper
provides, is the formal argument for \emph{why the hypothesis applies at
all}: why the transformer forward pass qualifies as an N-D-C
instantiation in the first place, such that the phase-boundary
predictions should hold.

This paper makes three contributions. First, it proves that the
transformer forward pass satisfies the formal definition of an N-D-C
recursion given in
\href{https://github.com/baardev/truevalue/raw/main/docnav/Research/papers/1_recursive-tholonic-five-constants/1_recursive-tholonic-five-constants.pdf}{Papers
1} and
\href{https://github.com/baardev/truevalue/raw/main/docnav/Research/papers/3_minimal-recursive-triadic-framework/3_minimal-recursive-triadic-framework.pdf}{3}:
the three-role structure is instantiated at every layer, not merely
analogized to one. Second, it provides a complete component-level role
assignment table, mapping each transformer component to its tholonic
role and giving the structural justification for each assignment. Third,
it derives the primary structural predictions of the framework (phase
boundaries at the five constants, universal C-dominance, the virial
balance condition, and the optimal-depth formula) directly from the role
assignment, showing them to be consequences of the assignment rather
than empirical observations layered on top of it.
\href{https://github.com/baardev/truevalue/raw/main/docnav/Research/papers/10_tholonic-neural-architecture/10_tholonic-neural-architecture.pdf}{Paper
10}'s measurements are then interpreted as tests of this derived
structure, not post-hoc fits to data.

This paper contains no new empirical results. It is a theory paper whose
claims are structural: if the role assignment holds, the predictions
must follow. The empirical question of whether those predictions hold is
addressed in
\href{https://github.com/baardev/truevalue/raw/main/docnav/Research/papers/10_tholonic-neural-architecture/10_tholonic-neural-architecture.pdf}{Paper
10}.

\begin{center}\rule{0.5\linewidth}{0.5pt}\end{center}

\subsection{1. Introduction}\label{introduction}

\subsubsection{1.1 The logical gap}\label{the-logical-gap}

Three papers in the tholonic series bear directly on the question of
neural network structure.

\href{https://github.com/baardev/truevalue/raw/main/docnav/Research/papers/1_recursive-tholonic-five-constants/1_recursive-tholonic-five-constants.pdf}{Paper
1} (\emph{Emergence of Classical Constants from a Minimal Recursive
Triadic Framework}) proves that the five classical constants
\(\varphi\), \(\sqrt{2}\), \(\ln 2\), \(e\), and \(\pi/4\) emerge as the
fixed points or limiting values of a specific family of three-variable
recurrences. The recurrences are specified by the tholonic grammar: one
variable in the \(N\) role (the running state), one in the \(D\) role
(the bounding parameter), and one in the \(C\) role (the accumulating
parameter). The constants are not chosen: they emerge structurally as
the only limits consistent with the grammar.

\href{https://github.com/baardev/truevalue/raw/main/docnav/Research/papers/3_minimal-recursive-triadic-framework/3_minimal-recursive-triadic-framework.pdf}{Paper
3} (\emph{A Minimal Recursive Triadic Framework for Self-Similar
Hierarchical Systems}) proves that three is the minimum number of
functional roles required for non-trivial recursive self-organization.
Given the axioms of binary state space and simplex topology, any system
that organizes itself recursively must instantiate at least one \(N\)
state, one \(D\) bounding variable, and one \(C\) accumulating variable.
The roles are irreducible: no two-variable system can support the same
convergence behavior under the same independence conditions.

\href{https://github.com/baardev/truevalue/raw/main/docnav/Research/papers/10_tholonic-neural-architecture/10_tholonic-neural-architecture.pdf}{Paper
10} (\emph{Neural Networks as Tholonic Systems}) applies these results
empirically: it measures the internal dynamics of 20 transformer models
and finds that data-driven phase boundaries match a tholonic constant at
a 75.5\% rate, with role-consistent placement patterns and evidence of
universal C-dominance.

The logical structure, as it stands, is: \emph{the framework produces
the five constants}
(\href{https://github.com/baardev/truevalue/raw/main/docnav/Research/papers/1_recursive-tholonic-five-constants/1_recursive-tholonic-five-constants.pdf}{Paper
1}), and \emph{the five constants appear in transformers}
(\href{https://github.com/baardev/truevalue/raw/main/docnav/Research/papers/10_tholonic-neural-architecture/10_tholonic-neural-architecture.pdf}{Paper
10}). What is missing is the middle step: \emph{transformers instantiate
the framework}. Without that step, the match in
\href{https://github.com/baardev/truevalue/raw/main/docnav/Research/papers/10_tholonic-neural-architecture/10_tholonic-neural-architecture.pdf}{Paper
10} is empirically interesting but theoretically unexplained. A skeptic
could reasonably argue that five constants might appear at phase
boundaries in any complex system simply because those constants are
common in mathematics, and that no structural claim has been made about
why these specific constants should appear in these specific roles.

This paper closes that gap. It argues that the transformer forward pass
is not merely consistent with the N-D-C grammar but \emph{instantiates
it}: satisfying the formal definition of
\href{https://github.com/baardev/truevalue/raw/main/docnav/Research/papers/3_minimal-recursive-triadic-framework/3_minimal-recursive-triadic-framework.pdf}{Paper
3}'s triadic recursion at every layer. If that argument holds, the five
constants must appear at the phase boundaries as a mathematical
consequence, not as a pattern to be discovered empirically. The
role-specific placement predictions follow directly, as does the virial
balance condition, the C-dominance prediction, and the optimal-depth
formula.
\href{https://github.com/baardev/truevalue/raw/main/docnav/Research/papers/10_tholonic-neural-architecture/10_tholonic-neural-architecture.pdf}{Paper
10} then becomes a test of the role assignment, not an independent
discovery.

\subsubsection{1.2 The TVPCI parallel}\label{the-tvpci-parallel}

The appropriate model for this paper is
\href{https://github.com/baardev/truevalue/raw/main/docnav/Research/papers/2_supply-chain-transparency-tvpci/2_supply-chain-transparency-tvpci.pdf}{Paper
2} (\emph{Phase-Resolved Transparency Classification in Commodity Supply
Chains: TVPCI}). That paper is, structurally, a role-assignment paper:
it maps the N, D, and C roles to specific, observable entities in the
supply-chain domain (phase state, bounding evidence, corroboration
depth), derives a scoring formalism from that mapping, and notes
explicitly that ``the role assignments follow from the same structural
argument'' as the mathematical framework ``without importing any
mathematical results from the companion paper into the scoring rules.''
\href{https://github.com/baardev/truevalue/raw/main/docnav/Research/papers/2_supply-chain-transparency-tvpci/2_supply-chain-transparency-tvpci.pdf}{Paper
2} is not an empirical paper; it is the formal justification for why the
supply-chain application is valid and what predictions follow from it.

This paper does the same for the AI domain. The domain is different
(transformer components rather than supply-chain phases), the
derivations are different (phase boundary positions and virial balance
rather than transparency scores), but the logical structure is
identical: prove the mapping, derive the consequences, leave the
empirical tests to a separate paper.

\subsubsection{1.3 Organization}\label{organization}

Section 2 restates the N-D-C framework in the form needed for this
paper, citing
\href{https://github.com/baardev/truevalue/raw/main/docnav/Research/papers/1_recursive-tholonic-five-constants/1_recursive-tholonic-five-constants.pdf}{Papers
1} and
\href{https://github.com/baardev/truevalue/raw/main/docnav/Research/papers/3_minimal-recursive-triadic-framework/3_minimal-recursive-triadic-framework.pdf}{3}
for proofs. Section 3 argues that the transformer forward pass satisfies
the formal definition of an N-D-C recursion. Section 4 provides the
complete component-level role assignment. Section 5 derives the
structural predictions. Section 6 covers edge cases (MoE routing,
residual connections, positional encodings). Section 7 identifies the
falsifiability criteria: what observations would invalidate the role
assignment. Section 8 relates this paper to
\href{https://github.com/baardev/truevalue/raw/main/docnav/Research/papers/10_tholonic-neural-architecture/10_tholonic-neural-architecture.pdf}{Paper
10} and discusses the implications for architecture design.

\begin{center}\rule{0.5\linewidth}{0.5pt}\end{center}

\subsection{2. The N-D-C Framework}\label{the-n-d-c-framework}

\subsubsection{2.1 The tholonic triad}\label{the-tholonic-triad}

The formal definition of the tholonic triad is given in
\href{https://github.com/baardev/truevalue/raw/main/docnav/Research/papers/3_minimal-recursive-triadic-framework/3_minimal-recursive-triadic-framework.pdf}{Paper
3} (Definition 2.1). For this paper we use the following compact
statement.

\textbf{Definition 2.1} (Tholonic triad, from
\href{https://github.com/baardev/truevalue/raw/main/docnav/Research/papers/3_minimal-recursive-triadic-framework/3_minimal-recursive-triadic-framework.pdf}{Paper
3}). A \emph{tholonic triad} is an ordered triple \((N, D, C)\) of
non-negative reals with three functionally distinct roles:

\begin{itemize}
\tightlist
\item
  \(N\) (\emph{negotiation}): the running state; the quantity being
  iteratively refined, emerging from the interaction of the other two.
\item
  \(D\) (\emph{definition/limitation}): the constraining, bounding
  variable; what limits or specifies the state.
\item
  \(C\) (\emph{contribution/integration}): the accumulating,
  synthesizing variable; what generates and integrates into the state.
\end{itemize}

The three roles are irreducible:
\href{https://github.com/baardev/truevalue/raw/main/docnav/Research/papers/3_minimal-recursive-triadic-framework/3_minimal-recursive-triadic-framework.pdf}{Paper
3} (Lemma 4.2) proves that no two-variable system satisfies the same
convergence properties under functional independence conditions. The
minimum structure for a self-sustaining recursive system is exactly
three roles.

\subsubsection{2.2 The recursion and its fixed
points}\label{the-recursion-and-its-fixed-points}

A \emph{tholonic recursion} is a sequence of triples \((N_k, D_k, C_k)\)
governed by an update rule of the form

\[N_{k+1} = f(N_k, D_k, C_k), \qquad (D_{k+1}, C_{k+1}) = g(D_k, C_k),\]

where \(f\) maps the current state through the interaction of the
constraining and accumulating variables, and \(g\) updates the
parameters according to a traversal class (Advancing, Self-redefined, or
Fixed; see
\href{https://github.com/baardev/truevalue/raw/main/docnav/Research/papers/1_recursive-tholonic-five-constants/1_recursive-tholonic-five-constants.pdf}{Paper
1}, Section 4).

\href{https://github.com/baardev/truevalue/raw/main/docnav/Research/papers/1_recursive-tholonic-five-constants/1_recursive-tholonic-five-constants.pdf}{Paper
1} (Propositions 5.1 through 5.5) proves that the limits of five
distinct traversal classes are:

\[\lim_{k\to\infty} N_k \;\in\; \{\pi/4,\; \varphi,\; e,\; \sqrt{2},\; \ln 2\}\]

where each constant is the unique fixed point of its traversal class.
The constants are not chosen post-hoc; they are the only values
consistent with the three-role grammar under the respective update
rules. Each constant carries a structural interpretation:

\begin{itemize}
\tightlist
\item
  \(\varphi\): the self-similar scaling ratio; the fixed point of the
  proportional balance recursion between \(D\) and \(C\).
\item
  \(\sqrt{2}\): the scaling invariant at the \(D/C\) differentiation
  boundary; the first non-trivial scaling step in a binary-branching
  structure.
\item
  \(\ln 2\): the information-theoretic compression limit; the entropy of
  a uniform binary choice, which is the minimum loss at any output
  compression step.
\item
  \(e\): the continuous growth rate at the initial expansion; the base
  of the natural exponential, which governs the rate at which an
  accumulating variable grows from a unit seed.
\item
  \(\pi/4\): the rotational symmetry measure; the only branch requiring
  three numerically distinct seeds and external parameter injection,
  corresponding to the most structurally constrained (least symmetric)
  form of balance.
\end{itemize}

\subsubsection{2.3 Self-similar nesting}\label{self-similar-nesting}

\href{https://github.com/baardev/truevalue/raw/main/docnav/Research/papers/3_minimal-recursive-triadic-framework/3_minimal-recursive-triadic-framework.pdf}{Paper
3} (Section 5) establishes that any tholonic triad can be embedded as a
single element (\(N\), \(D\), or \(C\)) in a higher-level triad, without
altering the three-role grammar. This produces an unbounded hierarchy:
each level resolves to a new \(N\) state, which becomes the parent \(N\)
for the level above.

The recursive structure is specifically:

\[\text{Parent } N \;\longrightarrow\; (D,\, C) \;\longrightarrow\; \text{Child } N \;\longrightarrow\; (D',\, C') \;\longrightarrow \;\text{Grandchild } N \;\longrightarrow \cdots\]

This self-similar nesting is what distinguishes a \emph{tholonic
instantiation} from a system that merely has three components. A system
is a tholonic instantiation if and only if: (a) its components satisfy
the three-role assignment (one \(N\), one \(D\), one \(C\)), (b) the
roles interact according to the tholonic update rule, and (c) the
structure recurs at multiple scales without altering the grammar.

\subsubsection{2.4 Two structural consequences used in this
paper}\label{two-structural-consequences-used-in-this-paper}

Two consequences of the framework are used in the derivations of Section
5:

\textbf{Consequence A (Phase boundary positions).} In a tholonic system
operating across \(L\) levels, the transitions between recursion levels
are governed by the five constants as fractional depth positions.
Specifically, if the system's depth is normalized to the interval
\([0, 1]\), the expected boundary positions are:

\begin{itemize}
\tightlist
\item
  Embedding expansion boundary: \(e^{-1} \approx 0.368\) (early region)
\item
  Scaling boundary (\(\sqrt{2}\) zone): interval \([0.20, 0.55]\)
\item
  Mid-network equilibrium (\(\varphi\) zone): interval \([0.45, 0.80]\)
\item
  Output compression boundary (\(\ln 2\) zone): interval near
  \(1 - \ln 2 \approx 0.307\) from the end
\item
  Rotational balance (\(\pi/4\) zone): near \(0.785\) fractional depth
\end{itemize}

These are theoretical predictions, not empirically fitted values.

\textbf{Consequence B (Virial balance condition).} At structural
equilibrium, the constraining (\(D\)) and integrating (\(C\)) roles must
be in proportional balance. For the tholonic recursion, the equilibrium
condition is

\[\sigma_D \approx \tfrac{1}{2}\,\sigma_C,\]

where \(\sigma_D\) and \(\sigma_C\) are measures of the activation
magnitude (or computational weight) of the \(D\) and \(C\) components
respectively. The factor of \(\frac{1}{2}\) arises from the structure of
the recursion: in the balanced fixed-point case (\(\varphi\) branch),
\(D\) and \(C\) must maintain a ratio of \(1 : 2\) to sustain the
self-similar proportionality. A system where
\(\sigma_D \ll \frac{1}{2}\sigma_C\) is \emph{C-dominant}: its
integrative capacity substantially exceeds its definitional constraint.

\begin{center}\rule{0.5\linewidth}{0.5pt}\end{center}

\subsection{3. The Transformer Forward Pass as an N-D-C
Recursion}\label{the-transformer-forward-pass-as-an-n-d-c-recursion}

\subsubsection{3.1 Formal statement}\label{formal-statement}

The central claim of this paper is the following:

\textbf{Proposition 3.1} (Transformer forward pass as N-D-C recursion).
\emph{The forward pass of a standard transformer architecture satisfies
the formal definition of a tholonic recursion (Definition 2.1 and
\href{https://github.com/baardev/truevalue/raw/main/docnav/Research/papers/3_minimal-recursive-triadic-framework/3_minimal-recursive-triadic-framework.pdf}{Paper
3}, Section 4) at every layer. Specifically: (a) the hidden state
\(h_l\) at each layer \(l\) occupies the \(N\) role; (b) the layer
normalisation operation occupies the \(D\) role; (c) the combined
attention and MLP projection operations occupy the \(C\) role; and (d)
the recurrence
\(h_{l+1} = f(h_l, \text{LN}(h_l), \text{Attn+MLP}(h_l))\) satisfies the
tholonic update rule. The three roles are functionally distinct and
irreducible.}

The proof is constructive. We verify each condition of Definition 2.1
and
\href{https://github.com/baardev/truevalue/raw/main/docnav/Research/papers/3_minimal-recursive-triadic-framework/3_minimal-recursive-triadic-framework.pdf}{Paper
3} (Lemma 4.2) in turn.

\subsubsection{\texorpdfstring{3.2 The hidden state as the \(N\)
role}{3.2 The hidden state as the N role}}\label{the-hidden-state-as-the-n-role}

The \(N\) role requires: (i) a running state that is iteratively
refined; (ii) that it emerge from the interaction of the bounding and
accumulating roles; and (iii) that it carry the coherent representation
passed to the next iteration.

The hidden state \(h_l \in \mathbb{R}^{d_\text{model}}\) satisfies all
three conditions. It is the quantity that is refined at each layer
(condition i). Its value after layer \(l\) is a function of both the
normalised input (the \(D\) component: the bounded version of
\(h_{l-1}\)) and the attention/MLP output (the \(C\) component: the
accumulated new content). And it is what is passed to layer \(l+1\) as
the starting state for the next recursion cycle (condition iii).

Condition (ii) requires some care. In the standard pre-norm transformer
formulation, the update at each layer is

\[h_l = h_{l-1} + \text{Attn}(\text{LN}_1(h_{l-1})) + \text{MLP}(\text{LN}_2(h_{l-1} + \text{Attn}(\text{LN}_1(h_{l-1})))).\]

The new \(N\) state \(h_l\) is the sum of the previous \(N\) state
\(h_{l-1}\), a contribution from the attention block (which operates on
the \(D\)-normalized input), and a contribution from the MLP block
(which also operates on a \(D\)-normalized intermediate). The new state
is therefore a function of the interaction between the \(D\) and \(C\)
components, mediated by the residual connection. This satisfies
condition (ii).

\subsubsection{\texorpdfstring{3.3 Layer normalisation as the \(D\)
role}{3.3 Layer normalisation as the D role}}\label{layer-normalisation-as-the-d-role}

The \(D\) role requires: (i) that it constrain or bound the state; (ii)
that it define what the state is allowed to be before transformation;
and (iii) that it be functionally distinct from and independent of the
\(C\) role.

Layer normalisation (LN) and its variant RMSNorm satisfy all three
conditions. LN operates by computing

\[\text{LN}(h) = \gamma \cdot \frac{h - \mu}{\sigma} + \beta,\]

where \(\mu\) and \(\sigma\) are the mean and standard deviation of
\(h\), and \(\gamma, \beta\) are learnable scale and shift parameters.
This operation: forces the hidden state onto a standardised manifold
(condition i: bounding); projects the raw hidden state into a form that
the attention or MLP block can interpret stably (condition ii: defines
what the state is before transformation); and is computed independently
of and prior to the attention and MLP computations (condition iii:
functional independence).

The critical structural property is that LN is a
\emph{pre-transformation operation}: it specifies what \(h_{l-1}\) is
(normalises its distribution) before the \(C\) components act on it.
This is exactly what the \(D\) role does in the tholonic grammar. The
\(D\) variable does not produce new content; it bounds the domain within
which the \(C\) variable operates. LN does not generate new information;
it constrains the input representation to a defined statistical
manifold.

\subsubsection{\texorpdfstring{3.4 Attention and MLP as the \(C\)
role}{3.4 Attention and MLP as the C role}}\label{attention-and-mlp-as-the-c-role}

The \(C\) role requires: (i) that it accumulate or integrate content;
(ii) that it generate outputs that contribute to the new state; and
(iii) that it be computationally distinct from and substantially larger
in weight than the \(D\) component.

The attention mechanism and MLP projection together satisfy these
conditions. Attention accumulates contextual information across the
sequence by computing weighted sums of value vectors:

\[\text{Attn}(Q, K, V) = \text{softmax}\!\left(\frac{QK^\top}{\sqrt{d_k}}\right) V.\]

This is an integration operation: it synthesizes information from
multiple positions in the sequence into a single representation. The MLP
block then applies a point-wise non-linear transformation, further
integrating the representation through a learned expansion and
contraction:

\[\text{MLP}(x) = W_2\,\sigma(W_1 x + b_1) + b_2.\]

Both operations \emph{contribute} new content to the hidden state:
attention contributes relational information (what other tokens are
relevant to this one), and MLP contributes compositional information
(what abstract features follow from the current representation). Neither
operation constrains what the representation is; both expand what it
becomes. This is the \(C\) role.

Condition (iii) of the \(C\) role refers to its dominance in parameter
count relative to the \(D\) component. In a standard transformer, the LN
parameters (\(\gamma, \beta \in \mathbb{R}^{d_\text{model}}\)) account
for \(2d_\text{model}\) parameters per layer, while the attention and
MLP blocks account for approximately \(12d_\text{model}^2\) parameters
(four weight matrices for attention, two for MLP). For
\(d_\text{model} = 768\) (GPT-2 base), the ratio is
\(\frac{12 \times 768^2}{2 \times 768} = 3072 : 1\). This is not a
C-dominance problem at the design level; it is the natural size ratio
between a normalization operation and the transformation it constrains.
What the tholonic framework predicts is that this ratio, if uncorrected
by training dynamics, produces C-dominance at the level of
\emph{activation magnitudes}, not merely parameter counts:
\(\sigma_D \ll \frac{1}{2}\sigma_C\). Whether training corrects for this
is the empirical question of
\href{https://github.com/baardev/truevalue/raw/main/docnav/Research/papers/10_tholonic-neural-architecture/10_tholonic-neural-architecture.pdf}{Paper
10}.

\subsubsection{3.5 Functional independence of the three
roles}\label{functional-independence-of-the-three-roles}

\href{https://github.com/baardev/truevalue/raw/main/docnav/Research/papers/3_minimal-recursive-triadic-framework/3_minimal-recursive-triadic-framework.pdf}{Paper
3} (Lemma 4.2) requires that the \(D\) and \(C\) components be
\emph{functionally independent}: neither should be a deterministic
function of the other, and swapping their roles should produce a
different dynamical behavior. For the transformer:

\begin{itemize}
\tightlist
\item
  LN (D) and Attention+MLP (C) are computationally independent: LN is
  applied to the raw hidden state and outputs a normalised version;
  Attention+MLP is applied to the LN output and generates new content.
  They are different operations with different parameters, applied in a
  fixed sequential dependency (D before C).
\item
  The roles cannot be swapped: applying the attention mechanism before
  normalisation and the normalisation afterwards is a different
  architecture (post-norm), with different training stability
  properties. The ordering D then C is not arbitrary.
\item
  Neither is a function of the other: the LN parameters
  \((\gamma, \beta)\) do not depend on the attention/MLP weights, and
  vice versa. They are trained independently.
\end{itemize}

These conditions satisfy Lemma 4.2's functional independence
requirement, confirming that the three roles are genuinely irreducible
in the transformer architecture.

\subsubsection{3.6 Summary of the proof}\label{summary-of-the-proof}

The transformer forward pass satisfies all conditions of Definition 2.1
and
\href{https://github.com/baardev/truevalue/raw/main/docnav/Research/papers/3_minimal-recursive-triadic-framework/3_minimal-recursive-triadic-framework.pdf}{Paper
3}'s triadic recursion:

\begin{itemize}
\tightlist
\item
  Hidden state \(h_l\): the \(N\) role (running state, refined at each
  layer, passed to the next).
\item
  LayerNorm: the \(D\) role (bounds the state, defines what it is before
  transformation).
\item
  Attention + MLP: the \(C\) role (accumulates and integrates, generates
  new content).
\item
  The update rule \(h_l = h_{l-1} + C(\text{LN}(h_{l-1}))\) matches the
  tholonic recurrence \(N_{k+1} = f(N_k, D_k, C_k)\).
\item
  The three roles are functionally independent and irreducible.
\end{itemize}

The transformer forward pass is therefore a tholonic instantiation: not
a system that resembles the N-D-C grammar, but one that formally
satisfies it. This is the structural premise from which all predictions
in Section 5 follow.

\begin{center}\rule{0.5\linewidth}{0.5pt}\end{center}

\subsection{4. Component-Level Role
Assignment}\label{component-level-role-assignment}

The full role assignment, with structural justification, is given in
Table 1. This table is the practical output of Section 3: it maps every
standard transformer component to its tholonic role and explains why
each assignment is the only consistent one.

\textbf{Table 1. Tholonic role assignments for transformer components.}

\begin{longtable}[]{@{}
  >{\raggedright\arraybackslash}p{(\linewidth - 4\tabcolsep) * \real{0.3333}}
  >{\raggedright\arraybackslash}p{(\linewidth - 4\tabcolsep) * \real{0.3333}}
  >{\raggedright\arraybackslash}p{(\linewidth - 4\tabcolsep) * \real{0.3333}}@{}}
\toprule\noalign{}
\begin{minipage}[b]{\linewidth}\raggedright
Component
\end{minipage} & \begin{minipage}[b]{\linewidth}\raggedright
Role
\end{minipage} & \begin{minipage}[b]{\linewidth}\raggedright
Structural justification
\end{minipage} \\
\midrule\noalign{}
\endhead
\bottomrule\noalign{}
\endlastfoot
Hidden state \(h_l\) & \(N\) & The coherent representation at each
layer; both the product of the previous cycle and the starting condition
for the next. \\
LayerNorm / RMSNorm & \(D\) & Pre-transformation normalisation; bounds
the representation to a defined statistical manifold before any
integrative operation acts on it. Defines what the state is; generates
no new content. \\
Attention heads & \(C\) (primary) & Integrates contextual information
across the sequence; accumulates relational content. The fundamental
contribution operation of the transformer. \\
MLP projection & \(C\) (secondary) & Integrates the relational
representation through non-linear expansion and contraction; contributes
compositional content. Together with attention, constitutes the full
\(C\) operation at each layer. \\
Residual connection & Structural scaffold & Carries \(N\) forward across
the \(D\) and \(C\) operations; preserves the N-to-N' recursion by
ensuring the new state is formed as a modification of the old, not a
replacement of it. Not a role; the mechanism that implements the
recursion. \\
Embedding layer & Entry-level \(N\) & Instantiates the initial coherent
representation from the token sequence; the \(N_0\) of the forward-pass
recursion. \\
Unembedding / LM head & Exit-level \(C\) & Projects the final \(N\)
state to a probability distribution over the vocabulary; the final
contribution of the forward pass to the output. \\
Positional encoding & \(D\) (structural context) & Defines the
positional structure within which the attention mechanism operates;
contributes to the bounding of what attention is permitted to mean. In
RoPE and ALiBi implementations, this is a per-layer definitional
constraint on the attention computation. \\
Dropout (during training) & \(D\) (regularisation) & Stochastic
suppression of contributions; a constraint that prevents the \(C\)
components from accumulating spurious patterns. Applied to the \(C\)
output before residual addition. \\
\end{longtable}

\subsubsection{4.1 Assignment uniqueness}\label{assignment-uniqueness}

The assignments in Table 1 are not arbitrary. The N-D-C roles are
distinguishable by three criteria that
\href{https://github.com/baardev/truevalue/raw/main/docnav/Research/papers/3_minimal-recursive-triadic-framework/3_minimal-recursive-triadic-framework.pdf}{Paper
3} uses to prove Lemma 4.2: a component occupies the \(D\) role if it
bounds without generating, the \(C\) role if it generates without
bounding, and the \(N\) role if it is both the product and the starting
condition of the recursion. In the transformer, no component satisfies
two of these criteria simultaneously, which is why the assignment is
unique.

LayerNorm bounds (standardises the distribution) but does not generate
new information: it is pure \(D\). Attention generates (produces new
contextual representations) but does not bound: it is pure \(C\). The
hidden state is both product and starting condition: it is \(N\). There
is no ambiguity, and no component is left unassigned.

\subsubsection{4.2 Assignment stability across architecture
variants}\label{assignment-stability-across-architecture-variants}

The assignment in Table 1 applies without modification to post-norm
transformers (original Vaswani et al.~architecture), pre-norm
transformers (GPT-2 and most current architectures), and RMSNorm
variants (LLaMA, Gemma). The structural role of LN/RMSNorm as the \(D\)
operation does not depend on its position relative to the attention
block; it depends on its function: normalising the input representation
before transformation.

For architectures without explicit LayerNorm (some early convolutional
models), the \(D\) role is occupied by whatever component performs the
bounding function: weight initialisation constraints, batch
normalisation, or gradient clipping. The framework does not require that
the \(D\) component be LayerNorm specifically; it requires that there be
a component that performs the bounding function. The prediction is that
architectures with a weaker or absent \(D\) component will be more
C-dominant and will show less role-consistent phase structure.

\begin{center}\rule{0.5\linewidth}{0.5pt}\end{center}

\subsection{5. Structural Predictions Derived from the Role
Assignment}\label{structural-predictions-derived-from-the-role-assignment}

Given the role assignment of Section 4 and the tholonic framework of
Section 2, five structural predictions follow. These are not empirical
observations; they are consequences of the assignment. The derivations
are given below.

\subsubsection{5.1 Phase boundaries at the five
constants}\label{phase-boundaries-at-the-five-constants}

The transformer forward pass, as established in Section 3, is a tholonic
recursion that proceeds across \(L\) layers. The framework (Section 2.3)
predicts that the transitions between structurally distinct recursion
phases are governed by the five constants. For a network of depth \(L\),
the predicted boundary positions as fractional depths are:

\textbf{Embedding expansion boundary} (\(e\) zone): The initial
expansion from token embedding to contextual representation corresponds
to the \(e\) branch of the tholonic recurrence. This is the phase where
the \(C\) component's accumulation rate is governed by the natural
exponential: the representation grows from a sparse, symbolic encoding
to a distributed contextual one. Predicted fractional position:
\(1/e \approx 0.368\) from the input end.

\textbf{Early scaling boundary} (\(\sqrt{2}\) zone): The first interior
boundary marks the transition from the early expansion phase to the
stable representation phase. This is where the \(D\) and \(C\)
components reach their first structural differentiation. The
\(\sqrt{2}\) constant governs the scaling step at the boundary between
the two-variable and three-variable regimes in
\href{https://github.com/baardev/truevalue/raw/main/docnav/Research/papers/3_minimal-recursive-triadic-framework/3_minimal-recursive-triadic-framework.pdf}{Paper
3}'s simplex argument, and it is the scaling invariant for binary
differentiation. Predicted fractional interval: \([0.20, 0.55]\).

\textbf{Mid-network equilibrium} (\(\varphi\) zone): The golden ratio
marks the interior equilibrium zone, where \(D\) and \(C\) are in
proportional balance. This is the zone where the representation is
neither expanding (early layers) nor compressing (late layers) but
maintaining a stable, self-similar structure. \(\varphi\) is the unique
fixed point of the proportional balance recursion, and it governs the
depth range at which the representation achieves maximum coherence.
Predicted fractional interval: \([0.45, 0.80]\).

\textbf{Output compression boundary} (\(\ln 2\) zone): The final phase
transition is from the stable representation phase to the output
compression phase, where the representation is compressed to a
vocabulary probability distribution. The information-theoretic minimum
loss at any binary compression step is \(\ln 2\), which is therefore the
governing constant at the output end of the network. Predicted
fractional position: near \(1 - \ln 2 \approx 0.307\) from the output
end.

\textbf{Rotational balance} (\(\pi/4\) zone): The \(\pi/4\) boundary is
the most constrained: it requires three numerically distinct seeds and
external parameter injection. In the transformer, this corresponds to
the zone where the positional encoding constraint (\(D\) component) is
maximally active relative to the attention integration (\(C\)
component). This is the zone of highest rotational symmetry in the
attention pattern: the point at which attention heads are most uniformly
distributed across positions. Predicted fractional position: near
\(0.785\) normalized depth.

These five predictions are directional: each constant is predicted to
appear in a specific role-consistent zone, not merely at any boundary.
This role-consistency is what distinguishes the tholonic prediction from
a claim that five constants might appear anywhere at five boundaries.
The claim is that each constant appears in the zone where its structural
interpretation is active.

\subsubsection{5.2 Universal C-dominance}\label{universal-c-dominance}

The virial balance condition (Section 2.4, Consequence B) predicts that
any architecture where \(\sigma_D \ll \frac{1}{2}\sigma_C\) is
structurally C-dominant. The role assignment (Table 1) identifies
LayerNorm as \(D\) and Attention+MLP as \(C\).

Given the parameter-count ratio of approximately \(3072:1\) (Section
3.4), any architecture whose training process does not specifically and
actively correct for the \(D/C\) imbalance will converge to a state
where the \(D\) component's activation magnitude is far below the virial
target. This is a structural prediction, not an empirical observation:
unless training explicitly embeds a D/C balance constraint, all standard
transformer architectures should show universal C-dominance.

The prediction is particularly strong because no standard training
objective (cross-entropy loss, perplexity minimisation, RLHF reward
maximisation) contains any term that would push \(\sigma_D\) toward
\(\frac{1}{2}\sigma_C\). The training process is blind to virial
balance. Therefore, C-dominance is the default state of any architecture
trained on standard objectives, regardless of architecture family or
scale.

\subsubsection{5.3 The virial balance
regulariser}\label{the-virial-balance-regulariser}

The structural prediction of universal C-dominance implies a concrete
training intervention: adding a virial balance term to the training loss
that penalises deviation from \(\sigma_D = \frac{1}{2}\sigma_C\). The
form of this regulariser, derived directly from the structural
condition, is

\[\mathcal{L}_\text{virial} = \lambda \sum_{l=1}^{L} \left| \sigma_D^l - \tfrac{1}{2}\sigma_C^l \right|^2,\]

where \(\sigma_D^l\) and \(\sigma_C^l\) are the root-mean-square
activation magnitudes of the LayerNorm and Attention+MLP outputs at
layer \(l\), and \(\lambda\) is a regularisation coefficient.

This regulariser is not a heuristic add-on. It is the direct
training-time enforcement of the structural equilibrium condition. A
model trained with \(\mathcal{L}_\text{virial}\) is a model whose
training objective includes an explicit pressure toward N-D-C balance at
every layer. The framework predicts that such a model should converge
faster (less energy is wasted correcting structural imbalance at each
step), show more coherent phase structure (boundaries at the predicted
positions), and exhibit greater robustness to distribution shift (a
structurally balanced model does not depend on the training distribution
to maintain its internal coherence).

\subsubsection{5.4 Optimal-depth
prediction}\label{optimal-depth-prediction}

The self-similar nesting property of the tholonic framework (Section
2.3) implies a relationship between the initial and final information
states of the network and the number of layers required to achieve the
transition. The \(\varphi\) branch of the tholonic recursion converges
at a rate governed by the golden ratio: at each recursion step, the
representation approaches its limit by a factor of \(1/\varphi\). The
number of steps required to reduce the initial entropy \(H_0\) to the
final entropy \(H_L\) is therefore

\[L^* = \log_\varphi\!\left(\frac{H_0}{H_L}\right).\]

This is a prediction about the optimal depth of a transformer with fixed
width and fixed training data: the number of layers beyond which adding
more layers yields diminishing returns at a rate faster than
\(1/\varphi\) per layer. It is not a prediction about a specific
architecture; it is a prediction about the functional relationship
between the network's entropy reduction and its depth.

\begin{center}\rule{0.5\linewidth}{0.5pt}\end{center}

\subsection{6. Edge Cases: MoE Routing, Residuals, and Positional
Encodings}\label{edge-cases-moe-routing-residuals-and-positional-encodings}

\subsubsection{6.1 Mixture-of-Experts
routing}\label{mixture-of-experts-routing}

In Mixture-of-Experts (MoE) architectures (Mixtral, DeepSeek-V2,
DeepSeek-V3), the standard MLP at each layer is replaced by a set of
expert MLPs, and a router network selects which experts to activate for
each token. The question for the role assignment is: what role does the
router occupy?

The router computes a distribution over experts and selects the
top-\(k\) for each token. This is a \emph{bounding} operation: it
restricts the set of experts that may contribute to the current token's
representation. It does not generate content; it defines which
content-generating components are active. By the assignment criterion of
Section 4.1, the router occupies a \(D\) role within the MoE layer.

The structural consequence is significant. A standard transformer has
one LayerNorm per attention or MLP sub-block as its \(D\) component. An
MoE transformer has the same LayerNorm plus a router, which adds a
second \(D\)-type operation at each layer. The framework predicts that
MoE architectures should therefore show a stronger \(D\) contribution
than their dense counterparts, moving the virial ratio
\(\sigma_D / \sigma_C\) closer to the target of \(\frac{1}{2}\), even if
the expert MLPs individually remain C-dominant. This is a testable
prediction that distinguishes MoE from dense architectures within the
tholonic framework.

\subsubsection{6.2 Residual connections}\label{residual-connections}

Residual connections are not a role: they are the mechanism that
implements the \(N\)-to-\(N'\) recursion. Their function is to ensure
that the new hidden state is formed as \(h_l = h_{l-1} + \Delta_l\)
rather than as \(h_l = \Delta_l\), preserving the previous \(N\) state
across the \(D\) and \(C\) operations.

In tholonic terms, the residual connection is the structural scaffolding
that makes the transformer a recursion (each \(h_l\) is formed from the
previous \(h_{l-1}\)) rather than a composition (each \(h_l\) is formed
from the input \(h_0\) directly). Without residuals, the forward pass
would be a sequential function composition, not a recursive refinement,
and the \(N\) role would collapse: there would be no persistent state to
refine. With residuals, each layer produces a new \(N\) state that is
recognisably a development of the previous one, consistent with the
tholonic recursion's requirement that \(N\) be both product and starting
condition.

The prediction from this analysis is that architectures without residual
connections, or with very shallow residual stacking, will show weaker
tholonic phase structure: fewer role-consistent boundary placements and
less coherent virial balance.

\subsubsection{6.3 Positional encodings}\label{positional-encodings}

Positional encodings (absolute PE in the original transformer, relative
PE in T5 and DeBERTa, Rotary PE in LLaMA and Gemma, ALiBi in MPT) all
perform the same structural function: they constrain the attention
mechanism by encoding the positional relationship between tokens. This
is a definitional operation: it defines what position means within the
attention structure.

The assignment to the \(D\) role (Table 1) applies to all positional
encoding schemes, but with one important nuance. In the original fixed
absolute PE (added to the embedding, not per-layer), the positional
constraint is applied once at the input level. In RoPE and ALiBi, the
positional constraint is applied at every layer within the attention
computation. The per-layer application is structurally stronger: it
continuously constrains the attention mechanism at every recursion step
rather than only at the entry. The framework predicts that per-layer
positional encoding schemes should produce more coherent phase structure
(the \(\pi/4\) zone in particular, which requires external parameter
injection, is structurally supported by per-layer constraints), while
single-application schemes should show weaker \(\pi/4\) zone structure.

\begin{center}\rule{0.5\linewidth}{0.5pt}\end{center}

\subsection{7. Falsifiability Criteria}\label{falsifiability-criteria}

A structural role assignment is only scientifically useful if it can be
falsified. The following observations would invalidate the role
assignment of Section 4:

\textbf{Criterion 1 (Role inconsistency of constants).} If the five
constants were found at phase boundaries but in structurally
inconsistent roles (e.g., \(e\) at the output compression zone,
\(\ln 2\) at the embedding expansion zone), the role assignment would be
falsified. A match rate above the 67\% threshold with no
role-consistency pattern is consistent with a random coincidence
interpretation and inconsistent with the structural derivation of
Section 5.1.

\textbf{Criterion 2 (Absence of universal C-dominance).} If any
significant fraction of standard-objective-trained transformer
architectures showed virial ratios near \(\frac{1}{2}\) without
explicitly training with the virial regulariser, the prediction of
Section 5.2 would be falsified.

\textbf{Criterion 3 (Virial regulariser ineffectiveness).} If training
with \(\mathcal{L}_\text{virial}\) produced no improvement in
convergence speed, no correction of virial balance, and no improvement
in OOD robustness, the structural interpretation of C-dominance as a
training deficit would be falsified. The regulariser could still be
ineffective without falsifying the role assignment (it might be that the
regulariser's gradient is too small to shift the equilibrium), but
ineffectiveness under a well-tuned \(\lambda\) would be strong evidence
against the structural account.

\textbf{Criterion 4 (Depth independence of phase boundaries).} If phase
boundaries occurred at the same absolute layer positions regardless of
network depth (e.g., always at layer 4, not at a constant fractional
depth), the fractal-scale independence required by the self-similar
nesting (Section 2.3) would be violated.

\textbf{Criterion 5 (MoE indistinguishability).} If MoE architectures
showed identical virial ratios to their dense counterparts of the same
parameter count, the prediction of Section 6.1 (that the router adds a
\(D\) component and shifts the virial ratio) would be falsified.

These criteria are specific and operationalizable.
\href{https://github.com/baardev/truevalue/raw/main/docnav/Research/papers/10_tholonic-neural-architecture/10_tholonic-neural-architecture.pdf}{Paper
10} and its future extensions provide evidence on Criteria 1 through 3.
Criteria 4 and 5 remain open experimental questions.

\begin{center}\rule{0.5\linewidth}{0.5pt}\end{center}

\subsection{\texorpdfstring{8. Relationship to
\href{https://github.com/baardev/truevalue/raw/main/docnav/Research/papers/10_tholonic-neural-architecture/10_tholonic-neural-architecture.pdf}{Paper
10}}{8. Relationship to Paper 10}}\label{relationship-to-paper-10}

\href{https://github.com/baardev/truevalue/raw/main/docnav/Research/papers/10_tholonic-neural-architecture/10_tholonic-neural-architecture.pdf}{Paper
10} (\emph{Neural Networks as Tholonic Systems}) is the empirical
counterpart of this paper. Its measurements are tests of the predictions
derived here. The relationship is asymmetric:

\begin{itemize}
\tightlist
\item
  This paper produces predictions from structure;
  \href{https://github.com/baardev/truevalue/raw/main/docnav/Research/papers/10_tholonic-neural-architecture/10_tholonic-neural-architecture.pdf}{Paper
  10} tests those predictions.
\item
  This paper cannot be falsified by
  \href{https://github.com/baardev/truevalue/raw/main/docnav/Research/papers/10_tholonic-neural-architecture/10_tholonic-neural-architecture.pdf}{Paper
  10}'s results alone (a failed empirical test might mean the
  measurements are wrong, not the assignment); but
  \href{https://github.com/baardev/truevalue/raw/main/docnav/Research/papers/10_tholonic-neural-architecture/10_tholonic-neural-architecture.pdf}{Paper
  10}'s results constrain the plausibility of the assignment.
\item
  If
  \href{https://github.com/baardev/truevalue/raw/main/docnav/Research/papers/10_tholonic-neural-architecture/10_tholonic-neural-architecture.pdf}{Paper
  10} fails its pre-registered criteria, the role assignment of Section
  4 should be reconsidered before revising the framework of
  \href{https://github.com/baardev/truevalue/raw/main/docnav/Research/papers/1_recursive-tholonic-five-constants/1_recursive-tholonic-five-constants.pdf}{Papers
  1} and
  \href{https://github.com/baardev/truevalue/raw/main/docnav/Research/papers/3_minimal-recursive-triadic-framework/3_minimal-recursive-triadic-framework.pdf}{3}.
\end{itemize}

The evidentially strongest result in
\href{https://github.com/baardev/truevalue/raw/main/docnav/Research/papers/10_tholonic-neural-architecture/10_tholonic-neural-architecture.pdf}{Paper
10} is not the headline 75.5\% match rate but the pattern of
role-consistent failures: the displacement of \(\varphi\) and
\(\sqrt{2}\) into structurally predicted directions in architectures
with known post-training alignment procedures. This pattern is predicted
by the derivation of Section 5.1 (post-training RLHF adds \(C\) content
without adding \(D\) constraint, displacing the interior equilibrium
zones toward the output) and is therefore more discriminating than a
simple boundary match rate.

The activation steering experiment
(\href{https://github.com/baardev/truevalue/raw/main/docnav/Research/papers/10_tholonic-neural-architecture/10_tholonic-neural-architecture.pdf}{Paper
10}, Section 8.2) is particularly relevant to this paper: the finding
that the \(\varphi\) zone has a shorter perturbation half-life than
surrounding zones in 100\% of tested models is a direct test of the N
role claim (Section 3.2) that the mid-network hidden state is the zone
of maximum coherence. A shorter perturbation half-life means the zone
more actively returns to its previous state after a perturbation, which
is the operational signature of an \(N\) attractor: a stable
configuration that restores itself. This is not a claim about boundary
positions; it is a claim about the functional character of the \(N\)
zone as a self-restoring state. The 100\% pass rate on this test is the
strongest single piece of evidence that the role assignment is
structurally correct.

\begin{center}\rule{0.5\linewidth}{0.5pt}\end{center}

\subsection{9. Discussion}\label{discussion}

\subsubsection{9.1 Implications for architecture
design}\label{implications-for-architecture-design}

The structural derivations of Section 5 suggest three concrete
directions for architecture design:

\textbf{Explicit D-component strengthening.} The universal C-dominance
prediction implies that any architectural change that increases the
relative weight of the \(D\) components (LayerNorm, positional
constraints, routing mechanisms) without proportionally increasing the
\(C\) components should move the architecture toward virial balance and
produce more coherent phase structure. Architectural choices that have
this effect include: increasing the number of normalization layers,
using per-layer rather than input-only positional encoding, and adding
routing constraints in MoE configurations.

\textbf{Virial balance as a training objective.} The virial regulariser
(Section 5.3) is the most direct architectural intervention: it embeds
the structural balance condition as a training signal. The framework
predicts that a model trained with this regulariser will achieve virial
balance by the end of training, regardless of its architectural starting
configuration. This is a softer intervention than architectural
modification: it adjusts the training dynamics rather than the forward
pass structure.

\textbf{Depth optimisation via the tholonic formula.} The optimal-depth
prediction (Section 5.4) provides a principled basis for choosing
network depth given a target entropy reduction. Current depth choices
are largely empirical (scale the model until performance saturates, then
add more data). The formula \(L^* = \log_\varphi(H_0/H_L)\) provides a
structural basis for depth selection that does not require exhaustive
scaling runs.

\subsubsection{9.2 Relationship to
alignment}\label{relationship-to-alignment}

Section 5.2 establishes that universal C-dominance is a structural
prediction of the role assignment, not an empirical curiosity. This has
implications for AI alignment that are worth stating explicitly.

Current alignment procedures (RLHF, Constitutional AI, RLAIF, DPO) all
share a structural property: they add preference information to the
model after the fact, via reward signals or preference data that modify
the \(C\) components (the attention and MLP weights) without
correspondingly strengthening the \(D\) components. This is not a design
flaw; it is a consequence of how these procedures work. Preference
learning operates by shaping what the model outputs (\(C\)
modification), not by strengthening what the model constrains itself to
be (\(D\) modification).

The tholonic framework predicts that this structural asymmetry is the
root cause of the characteristic failure modes of post-trained models:
specification gaming (the model satisfies the letter of the reward
signal while violating its intent, because the \(D\) constraint is too
weak to enforce the intent), distributional fragility (the model's
behavior degrades under distribution shift because its internal
coherence is maintained by training-distribution \(C\) patterns, not by
structural \(D\) constraints), and reward hacking (the model finds
\(C\)-side loopholes in the reward function because the \(D\)-side
constraints do not rule them out structurally).

The virial regulariser addresses this at the architectural level. By
training the model with an explicit balance constraint, the \(D\)
component is strengthened relative to the \(C\) component throughout
training, not just at the post-training stage. The structural prediction
is that a virial-balanced model is less susceptible to these failure
modes, not because it has been more thoroughly aligned post-hoc, but
because its internal constraint structure is intrinsically more robust.

\subsubsection{9.3 Scope}\label{scope}

This paper establishes a role assignment and derives structural
predictions. It does not claim that the tholonic framework explains
everything about transformer behavior, that the virial regulariser
solves the alignment problem, or that the five constants are the only
structurally significant features of transformer dynamics. The claims
are narrower: if the role assignment of Section 4 is correct, the
predictions of Section 5 must follow.
\href{https://github.com/baardev/truevalue/raw/main/docnav/Research/papers/10_tholonic-neural-architecture/10_tholonic-neural-architecture.pdf}{Paper
10} tests those predictions. The results are preliminary and all claims
should be understood as structurally motivated hypotheses requiring
independent verification.

\begin{center}\rule{0.5\linewidth}{0.5pt}\end{center}

\subsection{10. Conclusion}\label{conclusion}

This paper provides the missing middle step in the tholonic argument
about neural networks.
\href{https://github.com/baardev/truevalue/raw/main/docnav/Research/papers/1_recursive-tholonic-five-constants/1_recursive-tholonic-five-constants.pdf}{Paper
1} establishes that five classical constants emerge from the N-D-C
recurrence.
\href{https://github.com/baardev/truevalue/raw/main/docnav/Research/papers/3_minimal-recursive-triadic-framework/3_minimal-recursive-triadic-framework.pdf}{Paper
3} establishes that three roles are the minimum for recursive
self-organization.
\href{https://github.com/baardev/truevalue/raw/main/docnav/Research/papers/10_tholonic-neural-architecture/10_tholonic-neural-architecture.pdf}{Paper
10} measures phase boundaries empirically. What was missing was the
formal argument that transformer architectures instantiate the N-D-C
grammar, and that the empirical predictions therefore must follow
structurally rather than coincidentally.

The central result (Proposition 3.1) is that the transformer forward
pass satisfies all conditions of the formal tholonic recursion: hidden
states are \(N\), LayerNorm is \(D\), Attention+MLP is \(C\), and the
three roles are functionally independent and irreducible. Given this
assignment, five structural predictions follow directly: phase
boundaries at the five constants in role-consistent positions, universal
C-dominance in standard-objective-trained architectures, the form of the
virial balance regulariser, and the optimal-depth formula. These are not
post-hoc descriptions of observed patterns; they are derivations from
the role assignment that
\href{https://github.com/baardev/truevalue/raw/main/docnav/Research/papers/10_tholonic-neural-architecture/10_tholonic-neural-architecture.pdf}{Paper
10} then tests empirically.

The logical chain is now complete. The mathematical foundation
(\href{https://github.com/baardev/truevalue/raw/main/docnav/Research/papers/1_recursive-tholonic-five-constants/1_recursive-tholonic-five-constants.pdf}{Papers
1} and
\href{https://github.com/baardev/truevalue/raw/main/docnav/Research/papers/3_minimal-recursive-triadic-framework/3_minimal-recursive-triadic-framework.pdf}{3}),
the formal domain mapping (this paper), and the empirical test
(\href{https://github.com/baardev/truevalue/raw/main/docnav/Research/papers/10_tholonic-neural-architecture/10_tholonic-neural-architecture.pdf}{Paper
10}) together form a coherent, falsifiable structural account of why
transformer architectures organize their internal dynamics the way they
do, and what architectural changes would make them more stable, more
interpretable, and more robustly aligned.

\begin{center}\rule{0.5\linewidth}{0.5pt}\end{center}

\subsection{References}\label{references}

{[}Mil26a{]} Milton, J. W. \emph{Emergence of Classical Constants from a
Minimal Recursive Triadic Framework.} Clarity Coalition, paper 1 in this
series, 2026.

\url{https://github.com/baardev/truevalue/raw/main/docnav/Research/papers/1_recursive-tholonic-five-constants/1_recursive-tholonic-five-constants.pdf}

{[}Mil26b{]} Milton, J. W. \emph{Phase-Resolved Transparency
Classification in Commodity Supply Chains: A Structural Triadic Scoring
Framework (TVPCI).} Clarity Coalition, paper 2 in this series, 2026.

\url{https://github.com/baardev/truevalue/raw/main/docnav/Research/papers/2_supply-chain-transparency-tvpci/2_supply-chain-transparency-tvpci.pdf}

{[}Mil26c{]} Milton, J. W. \emph{A Minimal Recursive Triadic Framework
for Self-Similar Hierarchical Systems.} Clarity Coalition, paper 3 in
this series, 2026.

\url{https://github.com/baardev/truevalue/raw/main/docnav/Research/papers/3_minimal-recursive-triadic-framework/3_minimal-recursive-triadic-framework.pdf}

{[}Mil26d{]} Milton, J. W. \emph{Neural Networks as Tholonic Systems: A
Structural Framework for Architecture, Scaling, and
Alignment-by-Design.} Clarity Coalition, paper 10 in this series, 2026.

\url{https://github.com/baardev/truevalue/raw/main/docnav/Research/papers/10_tholonic-neural-architecture/10_tholonic-neural-architecture.pdf}

{[}VSP17{]} Vaswani, A., Shazeer, N., Parmar, N., Uszkoreit, J., Jones,
L., Gomez, A. N., Kaiser, L., Polosukhin, I. \emph{Attention Is All You
Need.} NeurIPS 2017.

{[}RNS+18{]} Radford, A., Narasimhan, K., Salimans, T., Sutskever, I.
\emph{Improving Language Understanding by Generative Pre-Training.}
OpenAI Technical Report, 2018.

{[}BMRS20{]} Brown, T., Mann, B., Ryder, N., et al.~\emph{Language
Models are Few-Shot Learners.} NeurIPS 2020.

{[}BA16{]} Ba, J., Kiros, J., Hinton, G. \emph{Layer Normalization.}
arXiv:1607.06450, 2016.

{[}ZSA+22{]} Zhang, B., Sennrich, R. \emph{Root Mean Square Layer
Normalization.} NeurIPS 2019.

{[}SLP+22{]} Shazeer, N. \emph{Mixture of Experts.} Various; see Fedus,
W., Zoph, B., Shazeer, N. \emph{Switch Transformers.} JMLR, 2022.

{[}JGS+21{]} Jiang, A. Q., et al.~\emph{Mistral 7B.} arXiv:2310.06825,
2023.

{[}Sha48{]} Shannon, C. E. \emph{A Mathematical Theory of
Communication.} Bell System Technical Journal, 27(3):379--423, 1948.

\begin{center}\rule{0.5\linewidth}{0.5pt}\end{center}

\emph{End of paper.}

\end{document}
